\section{Strict walks are NP-complete}\label{sec:strict}

A \emph{strict} walk has strictly increasing time labels and uses at most one edge of each snapshot. Let \SFIX{} and \SFREE{} denote the strict variants of \FIX{} and \FREE{}. The case with prescribed ends was left open by Deligkas et al.~\cite[p.~26504]{deligkas2025}, beyond the special cases of Section~\ref{sec:related}.

\subsection{Realizable orientations}

A strict walk that arrives at a vertex at time $t$ cannot leave it at time $t$.

\begin{lemma}\label{lem:strict}
An orientation is realizable by strict walks that start according to $a$ if and only if $\out_t(v)\le c_t(v)$ for every vertex $v$ and every time $t$. For \SFIX{} the final number of walks at $v$ must in addition equal $z_v$ for every $v$.
\end{lemma}

\begin{proof}
\emph{Necessity.} A strict walk uses at most one edge at time $t$. The $\out_t(v)$ arcs that leave $v$ at time $t$ are therefore traversed by pairwise distinct walks, and each of these walks is located at $v$ just before $t$. \emph{Sufficiency.} Process the times in increasing order. At time $t$, assign to the arcs that leave $v$ pairwise distinct walks located at $v$, which exist because $\out_t(v)\le c_t(v)$. Each of these walks traverses its arc and all other walks wait. The number of walks at every vertex after $t$ equals $c_{t^+}$, so the induction continues.
\end{proof}

Compared with Theorem~\ref{thm:char}, condition~(i) is replaced by the stronger condition $\out_t(v)\le c_t(v)$, and condition~(ii) disappears, since every edge is traversed by a walk that was at its tail before. Thus an orientation satisfying the strict condition has no dead even component. In the notation of Section~\ref{sec:char} the new condition reads $2B_{v,j-1}+b_{v,j}\le a_v+G_{v,j-1}$. The different weights on current and earlier departures prevent expressing it as a capacity constraint of the chain network.

\subsection{The reduction}

We reduce from satisfiability with clauses of two or three literals in which every variable occurs at most three times, which is NP-complete by Tovey~\cite{tovey1984}. We may assume that the literals of a clause belong to distinct variables (discard tautological clauses, merge repeated literals, propagate unit clauses), and that every literal occurs at most twice: a literal that occurs three times is pure, because its variable occurs at most three times and so the opposite literal does not occur, hence we set it to true and delete its clauses. This creates no new occurrences and no clause of size $1$. Let $\varphi$ be such a formula with variables $x_1,\dots,x_{n_\varphi}$ and clauses $C_1,\dots,C_{m_\varphi}$. Let $s_j\in\{2,3\}$ be the number of literals of $C_j$ and let $d(\ell)\le2$ be the number of clauses that contain the literal $\ell$. Since $\sum_js_j=\sum_\ell d(\ell)\le4n_\varphi$ and $s_j\ge2$, we have $m_\varphi\le2n_\varphi$. The temporal graph $G_\varphi$ has one vertex for each of the $2n_\varphi$ literals, one vertex for each clause, a \emph{hub} $h$, and vertices of degree one called \emph{leaves}. Its edges are the following (Fig.~\ref{fig:strict}).
\begin{itemize}
\item \emph{Time $1$.} One \emph{start leaf} is joined to every literal vertex, $s_j-1$ start leaves are joined to $C_j$, and $4n_\varphi$ start leaves are joined to $h$.
\item \emph{Time $2$.} The edge $\{x_i,\lnot x_i\}$ for every variable $x_i$.
\item \emph{Time $3$.} The edge $\{C_j,\ell\}$ for every clause $C_j$ and every literal $\ell$ of $C_j$.
\item \emph{Times $4$ and $5$.} One edge $\{\ell,h\}$ at time $4$ and one at time $5$, for every literal vertex $\ell$.
\item \emph{Time $6$.} $d(\ell)$ \emph{end leaves} are joined to every literal vertex $\ell$, one end leaf is joined to every $C_j$, and $6n_\varphi-2m_\varphi$ end leaves are joined to $h$.
\end{itemize}
We set $a_v=1$ at every start leaf, $z_v=1$ at every end leaf, and $a_v=z_v=0$ otherwise. There are $6n_\varphi+\sum_js_j-m_\varphi$ start leaves and as many end leaves. The graph $G_\varphi$ has no parallel edges within a snapshot, and its size is linear in the size of $\varphi$.

In Fig.~\ref{fig:strict}, edge labels are times and numbers beside leaves are multiplicities. White leaves are start leaves and shaded leaves are end leaves. Other clause incidences and the end leaves of $\lnot x_i$ are omitted.

\begin{figure}[!htbp]
\centering
\includegraphics{Fig8}
\caption{The strict-walk reduction around a variable and an incident clause.}
\label{fig:strict}
\end{figure}

\begin{theorem}\label{thm:strict}
\SFIX{} is NP-complete. It remains NP-complete when the temporal graph has six time labels and no parallel edges within a snapshot, and when every vertex with $a_v+z_v>0$ is a leaf with $a_v+z_v=1$, so that every walk of every solution is non-empty. \SFREE{} is NP-complete even for temporal graphs with two time labels, and with three time labels if empty walks are not allowed.
\end{theorem}

\begin{proof}
Both problems are in NP: an orientation is a certificate checked in polynomial time by Lemma~\ref{lem:strict}, and when empty walks are not allowed the certificate also lists the at most $m$ non-empty walks. We show that $\varphi$ is satisfiable if and only if $G_\varphi$ with the vectors $a$ and $z$ above is a feasible instance of \SFIX{}.

\emph{A solution yields a satisfying assignment.} Take the orientation of a solution. Before time $1$ only the start leaves hold walks, so the condition $\out_1(v)\le c_1(v)$ of Lemma~\ref{lem:strict} forces every edge at time $1$ to be directed away from its leaf, since its other end holds no walk. Before time $2$ every literal vertex therefore holds one walk, $C_j$ holds $s_j-1$ walks and $h$ holds $4n_\varphi$ walks. The edge $\{x_i,\lnot x_i\}$ at time $2$ moves one walk, so before time $3$ one of the two vertices $x_i$ and $\lnot x_i$ holds two walks and the other holds none. Call a literal \emph{true} if its vertex holds two walks. Exactly one of $x_i$ and $\lnot x_i$ is true, so this is a truth assignment. At time $3$ the vertex $C_j$ has $s_j$ edges but holds only $s_j-1$ walks. By Lemma~\ref{lem:strict} at least one of these edges is directed from a literal vertex $\ell$ to $C_j$, and $\ell$ holds a walk before time $3$. Hence $\ell$ is true and $C_j$ is satisfied.

\emph{A satisfying assignment yields a solution.} We give an orientation and verify Lemma~\ref{lem:strict}. Direct the edges at time $1$ away from the leaves and the edge $\{x_i,\lnot x_i\}$ from the false literal to the true literal. For every clause $C_j$ choose one true literal $\ell_j$ of $C_j$, direct $\{C_j,\ell_j\}$ from $\ell_j$ to $C_j$, and direct the other $s_j-1$ edges of $C_j$ away from $C_j$. Let $\theta_\ell\le d(\ell)\le2$ be the number of clauses that choose a true literal $\ell$. A true literal holds two walks, which suffice for its at most two outgoing edges at time $3$, and a false literal has none. After time $3$ the vertex $C_j$ holds one walk, a false literal $\ell$ holds $d(\ell)$ walks, a true literal $\ell$ holds $d(\ell)+2-2\theta_\ell$ walks, and $h$ holds $4n_\varphi$ walks. The edges at times $4$ and $5$ bring every literal vertex $\ell$ to exactly $d(\ell)$ walks.
\begin{itemize}
\item If $\ell$ is true and $\theta_\ell=0$, direct both edges from $\ell$ to $h$. The vertex $\ell$ holds $d(\ell)+2\ge2$ walks before time $4$.
\item If $\ell$ is false, or true with $\theta_\ell=1$, direct the edge at time $4$ from $h$ to $\ell$ and the edge at time $5$ from $\ell$ to $h$. The vertex $\ell$ holds at least one walk before time $5$.
\item If $\ell$ is true and $\theta_\ell=2$, direct both edges from $h$ to $\ell$.
\end{itemize}
At each of the times $4$ and $5$ at most $2n_\varphi$ edges leave $h$, and $h$ holds $4n_\varphi$ walks before time $4$ and at least $2n_\varphi$ walks before time $5$. The number of walks at $h$ changes in total by $\sum_{\ell\text{ true}}(2-2\theta_\ell)=2n_\varphi-2m_\varphi$, because there are $n_\varphi$ true literals and $\sum_\ell\theta_\ell=m_\varphi$. Hence $h$ holds $6n_\varphi-2m_\varphi\ge0$ walks before time $6$ (using $m_\varphi\le2n_\varphi$). Before time $6$ every vertex that is not a leaf thus holds exactly as many walks as it has end leaves, and the edges at time $6$ are directed towards the end leaves. The final numbers of walks are $z$.

\emph{Free end vertices.} Keep only the literal and clause vertices and the edges of times $2$ and $3$, and set $a_v=1$ at every literal vertex and $a_{C_j}=s_j-1$. Each vertex then holds the same number of walks before time $2$ as before, so the arguments for times $2$ and $3$ apply unchanged and this instance of \SFREE{} is feasible if and only if $\varphi$ is satisfiable. In a solution, a true literal that no clause chooses keeps an empty walk. If empty walks are not allowed, keep in addition the start leaves and the edges of time~$1$, with $a_v=1$ at every start leaf and $a_v=0$ elsewhere. Every walk then starts with an edge, and the same arguments apply with three time labels.
\end{proof}

The graph $G_\varphi$ is not proper and its snapshots have different numbers of edges, so neither Corollary~6.4 nor Theorem~6.5 of~\cite{deligkas2025} applies. The latter is exponential in $k$, which grows with $n_\varphi$ here. Theorem~5.4 of~\cite{deligkas2025} concerns the problem without prescribed start vertices and does not imply Theorem~\ref{thm:strict}. Whether \SFIX{} remains NP-complete when every vertex has at most two edges in every snapshot is left open.
